\section{Experiment Design}
\label{sec:study_design}
We characterize code quality differences across three analytical dimensions: stylistic violations, structural complexity, and security vulnerability patterns.

\subsection{RQ1: Structural and Stylistic Quality Analysis}

\noindent \textbf{RQ1:} \textit{What structural and stylistic quality differences exist between AI-generated and human-written code?}

RQ1 is examined across two analytical dimensions: stylistic violations and structural complexity.

For stylistic quality, we compute violation density (violations per file) as the primary comparison metric. Violations are further broken down by severity category, and the most frequent violation types per authorship class are identified to characterise predominant issue patterns. We stratify results by type discipline, comparing statically typed languages (C, C\texttt{++}, C\#, Java, Go) against dynamically typed languages (Python, JavaScript, PHP, Ruby) to examine whether static type systems moderate defect risk in AI-generated code. HTML is treated separately as a markup language and excluded from this typed-language comparison.

For structural complexity, we extract 19 metrics per file using the static analysis tool described in \S\ref{sec:static_analysis}, organised into two groups: 11 code stylometry metrics covering volume and lexical richness (Lines, Code Lines, Blank Lines, Comment Lines, Statements and their declarative and executable variants, Comment-to-Code Ratio, and Halstead Volume) and 8 complexity metrics covering control-flow and path density (Cyclomatic Complexity and its modified and strict variants, Essential Complexity, Maximum Nesting Depth, Paths, and Logarithmic Paths). These metrics enable direct comparison of structural properties between AI-generated and human-authored code and allow us to test whether traditional complexity-quality relationships established for human-written code generalise to AI-generated code.

Since code metric distributions are generally non-normal, all pairwise comparisons use nonparametric tests with effect sizes, applied consistently across both dimensions and repeated stratified by language to produce the type-discipline moderation analysis.

\subsection{RQ2: Security Vulnerability Profile Analysis}

\noindent \textbf{RQ2:} \textit{How do security vulnerability profiles differ between AI-generated and human-written code?}

To characterize security vulnerability profiles, we use Semgrep as the primary static analysis tool, applying its CWE-mapped rule sets uniformly across all ten languages. Vulnerability density, measured as findings per thousand lines of code, serves as the primary comparison metric since it accounts for file-length variation across the corpus. Each finding is mapped to its CWE identifier, allowing us to compare not just how many vulnerabilities appear in each origin but which vulnerability classes are disproportionately represented in AI-generated code. We separately examine the severity distribution of findings across info, warning, error, and critical levels, since a gap in raw density tells an incomplete story if the underlying findings differ substantially in exploitability risk.

Beyond the aggregate comparison, we report vulnerability density and CWE composition per language and then group results by type discipline to examine whether the same static-versus-dynamic moderation observed in RQ1 extends to the security dimension. We also compute a LOC-per-CWE measure, expressing on average how many lines of code are produced before a vulnerability appears, which gives a more intuitive view of security density differences between the two origins.


\subsection{RQ3: Open-Weight vs. Closed-Weight LLM Comparison}

\noindent \textbf{RQ3:} \textit{Do open-weight LLMs produce code with measurably different quality and security profiles compared to closed-weight LLMs when evaluated on the same benchmark dataset?}

A central challenge for RQ3 is the capability confound: if open-weight models produce lower quality code, the difference may reflect weaker model capability rather than weight openness itself. Since the models in CodeMirage are fixed, we adopt a stratified analysis rather than capability matching. The ten models are classified into open-weight (DeepSeek-V3, DeepSeek-R1, Llama-3.3-70B, Qwen-2.5-Coder-32B) and closed-weight (Claude-3.5-Haiku, GPT-4o-mini, GPT-o3-mini, Gemini-2.0-Flash, Gemini-2.0-Flash-Thinking, Gemini-2.0-Pro), with capability tiers assigned using publicly available HumanEval+ scores as a control variable. Reporting findings within each tier separately allows us to assess whether any observed gap persists after capability is roughly held constant.

The metric-based comparison reuses the measurement framework from RQ1 and RQ2, covering violation density, cyclomatic complexity, Halstead Volume, vulnerability density, CWE type distribution, and LOC-per-CWE, with model weight openness as the independent variable. We compare the aggregated open-weight group against the aggregated closed-weight group for a headline finding, then report per-model profiles to give practitioners a direct reference for each model in the benchmark. To complement this, we conduct a structural fingerprinting analysis using Tree Edit Distance (TED). We measure intra-group structural diversity by computing pairwise TED within each model family to examine whether open-weight and closed-weight models converge on similar tree structures. We also compute distances from each model family to the human baseline to see how closely each group aligns with human-written code in terms of structure. This structural analysis is independent of model capability and connects the AST-level findings from RQ1 directly to the model-level comparison in RQ3.
